\begin{figure}[htbp]
\centering
\begin{tikzpicture}[
  scale=0.95, transform shape,
  box/.style={draw, rounded corners, fill=green!80!white!10, align=center, text width=4.2cm, minimum height=.8cm, font=\small},
  boxaux/.style={draw, rounded corners, dashed, fill=green!20!cyan!15, align=center, text width=2.6cm, minimum height=1.1cm, font=\small},
  edg/.style={-{Latex[length=3mm]}, thick},
  edgd/.style={-{Latex[length=3mm]}, dashed, gray!80!black},
  every node/.style={font=\small}
]
\node[box] (S00) at (0,0) {00: Entorno};
\node[box] (S00b) at (0,-1.5) {00b: Modelos candidatos};
\node[boxaux] (S00c) at (5.5,-1.5) {00c\\Verif. literatura};
\node[box] (S01) at (0,-3) {01: Catálogo ESGF};
\node[box] (S02) at (-2.6,-4.5) {02: Descarga CMIP6};
\node[box] (S03) at (3.9,-4.5) {03: Descarga ERSSTv5};
\node[box] (S02b) at (-2.6,-6) {02b: Nodos ESGF alt.};
\node[box] (S02c) at (-2.6,-7.5) {02c: Copernicus CDS};
\node[box] (S04) at (0.5,-9) {04: Grilla común};
\node[box] (S05) at (0.5,-10.5) {05: Máscara océano-tierra};
\node[box] (S06) at (0.5,-12) {06: Control de calidad};
\node[box] (S07) at (0.5,-13.5) {07: Inventario final};
\node[box] (S08) at (0.5,-15) {Registro de modelos};
\node[box] (S09) at (0.5,-16.5) {Gráficos, reporte y paquete};
\draw[edg] (S00) -- (S00b);
\draw[edgd] (S00b) -- (S00c);
\draw[edg] (S00b) -- (S01);
\draw[edg] (S01) -- (S02);
\draw[edg] (S02) -- (S02b);
\draw[edg] (S02b) -- (S02c);
\draw[edgd] (S02b) to[bend right=60] (S01);
\draw[edgd] (S02c) to[bend right=80] (S01);
\draw[edg] (S02.west) -- ++(-1.5,0) |- (S04.west);
\draw[edg] (S03) -- (S04);
\draw[edg] (S04) -- (S05);
\draw[edg] (S05) -- (S06);
\draw[edg] (S06) -- (S07);
\draw[edg] (S07) -- (S08);
\draw[edg] (S08) -- (S09);
\end{tikzpicture}
\caption[Flujograma del \emph{pipeline} del Entregable~1.]{Flujograma del \emph{pipeline} del Entregable~1. Las flechas sólidas son transiciones de la secuencia automática de \texttt{run.sh}. Las punteadas corresponden al paso diagnóstico manual \texttt{00c} y a la actualización del catálogo del Paso~01 que hacen los pasos 02b y 02c cuando resuelven un modelo por una fuente alternativa. La cascada 02~$\to$~02b~$\to$~02c es automática dentro del Paso~02: cada fuente se intenta solo para lo que la anterior no resolvió. Las dos últimas cajas se ejecutan siempre, sin importar el rango de pasos pedido.}
\label{fig:flujograma}
\end{figure}
